\subsubsection{Extracción de datos}

La extracción de datos se diseñó para responder de forma directa la pregunta principal de investigación y las directrices D1--D3, manteniendo trazabilidad completa desde cada afirmación de síntesis hasta su evidencia en texto completo. El diseño del instrumento tomó como base el principio de Kitchenham de definir previamente qué datos son necesarios para contestar las preguntas de investigación y cómo deben codificarse para reducir sesgo de interpretación entre estudios.

Antes de la extracción definitiva se ejecutó un piloto sobre un subconjunto de artículos incluidos, con el objetivo de estabilizar definiciones operacionales, reducir ambigüedad semántica y verificar que cada campo aportara valor analítico real. Ese piloto permitió consolidar una matriz operativa estructurada en seis bloques: metadatos bibliográficos, datos y variables (D1), enfoque de lógica difusa (D2), plataforma y operación en tiempo real (D3), resultados de evaluación y trazabilidad metodológica. Como resultado de esta depuración, la matriz final quedó en 48 columnas, evitando campos redundantes y conservando únicamente variables útiles para síntesis transversal.

La ejecución de la extracción se realizó artículo por artículo y criterio por criterio, registrando tanto variables categóricas como descripciones técnicas y valores cuantitativos reportados. Para las variables categóricas se aplicó un catálogo controlado con vocabulario predefinido y reglas de normalización. Cuando una categoría no estaba cubierta por el catálogo, se usó el formato \texttt{Other:<texto>} para preservar precisión sin perder consistencia. El valor \texttt{NR} se reservó para información no reportada y \texttt{NA} para no aplicabilidad metodológica. Esta convención permitió separar ausencia de reporte, no aplicabilidad y variación terminológica genuina.

Toda celda analítica de D1, D2, D3 y evaluación se respaldó con evidencia de página en formato \texttt{p.X} o \texttt{p.X-p.Y}. La evidencia se registró de forma independiente por bloque (\texttt{Evidence\_D1\_pages}, \texttt{Evidence\_D2\_pages}, \texttt{Evidence\_D3\_pages}, \texttt{Evidence\_Eval\_pages}) para facilitar auditoría focalizada y revisión cruzada durante la síntesis. Esta estrategia evitó depender de notas narrativas no trazables y mejoró la reproducibilidad del proceso.

La consistencia de la extracción se verificó en dos niveles. Primero, se aplicaron controles estructurales y semánticos sobre la matriz completa: completitud de campos críticos, conformidad con catálogo, formato de evidencia y coherencia entre columnas relacionadas. Segundo, se ejecutó una auditoría de veracidad contra PDFs para confirmar que las páginas citadas efectivamente sustentaran contaminantes, diseño difuso, arquitectura/plataforma y métricas reportadas. Con estos controles, la matriz final de extracción quedó validada para los 19 estudios incluidos en texto completo.

La implementación operativa de esta fase se conserva en \texttt{articles\_data\_extraction.csv}, y su codificación controlada en \texttt{catalogo\_controlado\_extraccion\_v1.md}. En la versión de artículo, estos insumos deben citarse en el paquete reproducible consolidado (Zenodo, DOI único) junto con el checklist de calidad y el consolidado de resultados. La vista conceptual de la matriz se resume en la Tabla~\\ref{tab:matriz_extraccion_actualizada} y el detalle campo a campo se documenta en el Anexo~\\ref{sec:anexo_diccionario_extraccion}.
